\documentclass[../main.tex]{subfiles}

\begin{document}

\subsection{Narrative Style}
\label{subsec:narrative-style}

The narrative style of the Vietnam Financial Inclusion Dashboard 2025 is
\textbf{institutional, evidence-based, and policy-diagnostic}. The dashboard
does not use emotional appeals or individual anecdotes. Instead, it builds its
argument through measured indicators, structured comparisons, subgroup
breakdowns, and diagnostic gaps. This style is appropriate for a dashboard
based on weighted Global Findex microdata because its purpose is to support
policy interpretation rather than promotion or entertainment.

The overall tone can be summarized as:

\begin{center}
\textbf{Neutral evidence \(\rightarrow\) Diagnostic comparison \(\rightarrow\) Policy interpretation}
\end{center}

In this sense, the dashboard acts like an analytical policy brief. It does not
only report whether financial inclusion exists; it shows where inclusion is
strong, where it weakens, and which population groups or payment flows require
closer attention.

\noindent\textit{See Figures~\ref{fig:dashboard-overview},
\ref{fig:page2-readiness-overview}, \ref{fig:account-kpi-cards},
\ref{fig:national-inflows}, and \ref{fig:page5-resilience-overview}.}

\subsubsection{Evidence-Based and Diagnostic Tone}
\label{subsubsec:evidence-based-diagnostic-tone}

The dashboard lets quantitative gaps carry the argument. On Page 2, Online
Digital ID reaches 71.19\%, while Active Digital ID Verification is only
32.68\%. This 38.51 percentage-point gap supports the interpretation that
digital access does not automatically become active digital use. The Digital
Activity Profile reinforces the same neutral tone: Social Media activity reaches
94.46\%, while Gov activity is 15.63\% and Job activity is only 4.89\%. The
visual language remains factual, allowing the reader to infer the problem from
measurable contrasts.

The dashboard is also diagnostic because it does not stop at national-level
outcomes. Page 3 reports Banking Inclusion at 74.93\%, Financial Institution
Account ownership at 74.05\%, Digital Account ownership at 67.03\%, and Mobile
Money Account ownership at 46.69\%. However, the Exclusion Friction Log then
breaks remaining exclusion into barriers such as Not Enough Money, Too Far
Away, Too Expensive, Lack of Trust, and Lack of Documentation. This changes the
narrative from a descriptive statement about account access into a diagnostic
question about why exclusion persists.

\dashboardfigure{0.86}{fig-page2-digital-activity-profile.png}
{Digital activity profile: penetration depth among internet users.}
{fig:page2-digital-activity-profile}

\noindent\textit{See also Figures~\ref{fig:page2-active-digital-id},
\ref{fig:account-kpi-cards}, and \ref{fig:exclusion-friction-log}.}

\subsubsection{Comparative Policy Narrative}
\label{subsubsec:comparative-policy-narrative}

A major feature of the dashboard's narrative style is comparison. Page 4 shows
that Corporate/Formal Wages are mostly digitized, with 80.99\% delivered
directly to an account and 18.60\% cash-only. Agricultural Sales show the
opposite pattern: 73.42\% are cash-only, while only 25.79\% are delivered
digitally to an account. Government Subsidies/Transfers are also more
cash-oriented, with 56.96\% cash-only and 35.35\% digital direct-to-account.

The same comparative style appears in the urban--rural payment gap. Urban
Merchant Payments reach 65.79\%, compared with 50.30\% in rural areas. Urban
In-Store Digital Pay reaches 62.67\%, compared with 46.11\% in rural areas.
These paired comparisons turn the dashboard from a generic report about digital
payments into a policy narrative about uneven digitization across sectors and
geographies.

\noindent\textit{See Figures~\ref{fig:national-inflows} and
\ref{fig:page4-digital-merchant-payments}.}

\subsubsection{Policy-Oriented and Outcome-Oriented Ending}
\label{subsubsec:policy-outcome-ending}

The dashboard's narrative style is policy-oriented because each major gap
points toward an area for further investigation. The identity gap suggests a
need to look beyond ID availability toward user activation. The mobile money
gap suggests that account access differs by channel. The cash dependence of
Agricultural Sales suggests that digitization is weaker in some payment flows
than in formal wage channels. The rural payment gap suggests uneven merchant
acceptance and payment infrastructure.

The final page gives the story an outcome-oriented endpoint. Page 5 reports
that 73.05\% of respondents saved, while 45.98\% borrowed. More importantly,
the coping-channel table shows that the Poorest 20\% rely on Family \& Friends
at 86, while Formal Savings is only 30. In contrast, the Richest 20\% show
Formal Savings at 60 and Family \& Friends at 34. This ending prevents the
dashboard from treating financial inclusion as an end in itself. The final
question is whether people have the capacity to save, borrow safely, and
respond to financial shocks.

\noindent\textit{See Figures~\ref{fig:page5-resilience-overview} and
\ref{fig:page5-coping-income-group}.}

\subsubsection{Summary of Narrative Style}
\label{subsubsec:narrative-style-summary}

Table~\ref{tab:narrative-style-summary} summarizes how the dashboard's
narrative style is supported by selected visual evidence.

\begin{table}[htbp]
\centering
\small
\caption{Summary of narrative style and supporting evidence}
\label{tab:narrative-style-summary}
\begin{tabular}{p{0.24\textwidth}p{0.66\textwidth}}
\hline
\textbf{Narrative quality} & \textbf{Supporting evidence from the dashboard} \\
\hline
Institutional & Formal policy-dashboard language and weighted Global Findex indicators across the five dashboard pages. \\
\hline
Evidence-based & Online Digital ID = 71.19\%; Active Digital ID Verification = 32.68\%. \\
\hline
Diagnostic & The Exclusion Friction Log decomposes unbanked barriers on the Account Ownership page. \\
\hline
Comparative & Financial Institution Account = 74.05\% versus Mobile Money Account = 46.69\%; Urban Merchant Payments = 65.79\% versus Rural Merchant Payments = 50.30\%. \\
\hline
Policy-oriented & Agricultural Sales cash-only = 73.42\%; rural payment gaps and mobile money gaps point to areas for further policy investigation. \\
\hline
Outcome-oriented & \% Who Saved = 73.05\%, \% Who Borrowed = 45.98\%, followed by income-based coping-channel differences. \\
\hline
\end{tabular}
\end{table}

Overall, the dashboard uses a neutral and evidence-based narrative style to
identify structural barriers in Vietnam's financial inclusion landscape. This
style prepares the reader for the next analytical dimension: how interaction
and exploration allow users to move from national indicators into subgroup-level
policy investigation.

\end{document}
